\begin{frame}{Teoretická východiska}

\begin{enumerate}
\item Virtualizace a hypervisory
\begin{itemize}
 \item Typy virtualizace (plná, paravirtualizace, kontejnery), KVM/QEMU
\end{itemize}
\vspace{0.2cm}
\item Distribuované souborové systémy
\begin{itemize}
 \item Principy distribuovaného úložiště
\end{itemize}
\vspace{0.2cm}
\item Clusterová architektura / High Availability
\begin{itemize}
 \item Koncepty HA, failover, Pacemaker, RDMA/IPoIB
\end{itemize}
\note{
\textbf{Virtualizace a hypervisory}
\begin{itemize}
	\item Více OS na jednom HW (KVM/QEMU)
	\item V práci testuji platformy nad clusterem
\end{itemize}
\textbf{Distribuované souborové systémy}
\begin{itemize}
	\item Sdílení úložiště napříč uzly
	\item Ceph, GlusterFS, Lustre, BeeGFS, NFS
\end{itemize}
\textbf{InfiniBand / RDMA / IPoIB}
\begin{itemize}
	\item Vysokorychlostní propojení bez zátěže CPU
	\item Klíčové pro výběr OS (Rocky = plná podpora)
\end{itemize}
}
\end{enumerate}
\end{frame}

